\chapter{2013年安徽卷（文）}

\section{单选题}

\begin{problem}
设 \(t\) 是单位向量，若复数 \(a = \frac{1}{3 - \i} (a \in \R)\) 是纯虚数，则 \(a\) 的值为（\quad）

\choices
    {\(-3\)}
    {\(-1\)}
    {1}
    {3}
\end{problem}

\begin{problem}
已知 \(A = \{x \mid x + 1 > 0\}\)，\(B = \{-2, -1, 0, 1\}\)，则 \(\left(\complement_B A\right) \cap B =\)（\quad）

\choices
    {\(\{-2, -1\}\)}
    {\(\{-2\}\)}
    {\(\{-1, 0, 1\}\)}
    {\(\{0, 1\}\)}
\end{problem}

\begin{problem}
如图所示，程序框图 (算法流程图) 的输出结果是（\quad）

\choices
    {\( \frac{3}{4} \)}
    {\( \frac{1}{6} \)}
    {\( \frac{11}{12} \)}
    {\( \frac{25}{24} \)}
\end{problem}

\begin{problem}
“\((2x - 1)x = 0\)”是“\(x = 0\)”的（\quad）

\choices
    {充分不必要条件}
    {必要不充分条件}
    {充分必要条件}
    {既不充分也不必要条件}
\end{problem}

\begin{problem}
若某公司从五位大学毕业生甲、乙、丙、丁、戊中录用三人，这五人被录用的机会均等，则甲或乙被录用的概率为（\quad）

\choices
    {\( \frac{2}{3} \)}
    {\( \frac{2}{5} \)}
    {\( \frac{2}{5} \)}
    {\( \frac{9}{10} \)}
\end{problem}

\begin{problem}
直线 \(x + 2y - 5 + \sqrt{5} = 0\) 被圆 \(x^{2} + y^{2} - 2x - 4y = 0\) 截得的弦长为（\quad）

\choices
    {1}
    {2}
    {4}
    {\( 4\sqrt{6} \)}
\end{problem}

\begin{problem}
设 \(S_{n}\) 为等差数列 \(\{a_{n}\}\) 的前 \(n\) 项和，\(S_{8} = 4a_{8}, a_{7} = -2\)，则 \(a_{6} =\)（\quad）

\choices
    {\(-6\)}
    {\(-4\)}
    {\(-2\)}
    {2}
\end{problem}

\begin{problem}
函数 \( y = f(x) \) 的图象如图所示，在区间 \([a, b]\) 上可找到 \( n (n \geqslant 2) \) 个不同的数 \( x_1, x_2, \cdots, x_n \)，使得 \( \frac{f(x_1)}{x_1} = \frac{f(x_2)}{x_2} = \cdots = \frac{f(x_n)}{x_n} \)，则 \( n \) 的取值范围为（\quad）

\choices
    {[2, 3]}
    {[2, 3, 4]}
    {[3, 4]}
    {[3, 4, 5]}
\end{problem}

\begin{problem}
设 \(\triangle ABC\) 的内角 \(A, B, C\) 所对边的长分别为 \(a, b, c\)，若 \(b + c = 2a\)，\(3\sin A = 5\sin B\)，则角 \(C =\)（\quad）

\choices
    {\( \frac{\pi}{3} \)}
    {\( \frac{2\pi}{3} \)}
    {\( \frac{3\pi}{4} \)}
    {\(\frac{5\pi}{6}\)}
\end{problem}

\begin{problem}
已知函数 \( f(x) = x^3 + ax^2 + bx + c \) 有两个极值点 \( x_1, x_2 \)，若 \( f(x_1) = x_1 < x_2 \)，则关于 \( x \) 的方程 \( 3[f(x)]^2 + 2af(x) + b = 0 \) 的不同实根个数为（\quad）

\choices
    {3}
    {4}
    {5}
    {6}
\end{problem}

\section{填空题}

\begin{problem}
函数 \( y = \ln \left(1 + \frac{1}{x}\right) + \sqrt{1 - x^{2}} \) 的定义域为 \fillinblank{} \fillinblank{}.
\end{problem}

\begin{problem}
若非负变量 \(x, y\) 满足的约束条件 \(\left\{ \begin{array}{l} x - y \geqslant -1, \\ x + 2y \leqslant 4. \end{array} \right.\) 则 \(x + y\) 的最大值为 \fillinblank{} \fillinblank{}.
\end{problem}

\begin{problem}
若非零向量 \(a, b\) 满足 \(\|a\| = 3 \|b| = |a + 2b|\)，则 \(a\) 与 \(b\) 夹角的余弦值为 \fillinblank{} \fillinblank{}.
\end{problem}

\begin{problem}
定义在 \(\R\) 上的函数 \(f(x)\) 满足 \(f(x + 1) = 2f(x)\). 若当 \(0 \leqslant x \leqslant 1\) 时，\(f(x) = x(1 - x)\)，则当 \(-1 \leqslant x \leqslant 0\) 时，\(f(x) = \fillinblank{}\).
\end{problem}

\begin{problem}
如图，正方体 \(ABCD - A_{1}B_{1}C_{1}D_{1}\) 的棱长为 1, \(P\) 为 \(BC\) 中点，\(Q\) 为线段 \(CC_{1}\) 上的动点，过 \(A\)、\(P\)、\(Q\) 的平面截该正方体所得的截面记为 \(S\)，则下列命题正确的是 \fillinblank{}. (写出所有正确命题的编号)

\circled{1} 当 \(0 < CQ < \frac{1}{2}\) 时，\(S\) 为四边形; \circled{2} 当 \(CQ = \frac{1}{3}\) 时，\(S\) 为等腰梯形; \circled{3} 当 \(CQ = \frac{3}{4}\) 时，\(S\) 与 \(C_1D_1\) 交点 \(R\) 满足 \(C_1R = \frac{1}{3}\); \circled{4} 当 \(\frac{3}{4} < CQ < 1\) 时，\(S\) 为六边形; \circled{5} 当 \(CQ = 1\) 时，\(S\) 的面积为 \(\frac{\sqrt{6}}{2}\) \fillinblank{}.
\end{problem}

\begin{problem}
设函数 \( f(x) = \sin x + \sin \left(x + \frac{\pi}{3}\right) \). (1) 求 \( f(x) \) 的最小值，并求使 \( f(x) \) 取得最小值的 x 的集合; (2) 不画图，说明函数 \( y = f(x) \) 的图象可由 \( y = \sin x \) 的图象经过怎样的变化得到 \fillinblank{}.
\end{problem}

\begin{problem}
为调查甲，乙两校高二年级学生某次联考数学成绩情况. 用简单随机抽样，从这两校中各抽取 30 名高三年级学生，以他们的数学成绩 (百分制) 作为样本，样本数据的茎叶图如图: 甲 乙 7 4 5 3 3 2 5 3 3 8 5 5 4 3 3 3 1 0 0 6 0 0 1 1 2 2 3 3 5 8 6 6 2 2 1 1 0 0 7 0 0 2 2 2 3 3 6 6 9 7 5 4 4 2 8 1 1 5 5 8 2 0 9 0 (1) 若甲校高三年级每位学生被抽取的概率为 0.05，求甲校高三年级学生总人数，并估计甲校高三年级这次职教考数学成绩的及格率 (60 分及 60 分以上为及格); (2) 设甲、乙两校高三年级学生这次职教考平均成绩分别为 \( x_{1}, x_{2} \) ，估计 \( x_{1} = x_{2} \) 的值 \fillinblank{}.
\end{problem}

\begin{problem}
如图，因棱锥 \(P - ABCD\) 的底面 \(ABCD\) 是边长为 2 的菱形，\(\angle BAD = 60^{\circ}\). 已知 \(PB = PD = 2\)，\(PA = \sqrt{6}\). (1) 证明: \(PC \perp BD\); (2) 若 E 为 PA 的中点，求三棱锥 P - BCE 的体积 \fillinblank{}.
\end{problem}

\begin{problem}
设函数 \( f(x) = ax - (1 + \alpha^2)x^2 \)，其中 \( \alpha > 0 \)，区间 \( I = \{x \mid f(x) > 0\} \). (1) 求 I 的长度 (注: 区间 \( (\alpha, \beta) \) 的长度定义为 \( \beta - \alpha \) ); (2) 给定常数 \(k \in (0,1)\)，当 \(1 - k \leqslant a \leqslant 1 + k\) 时，求 \(I\) 长度的最小值 \fillinblank{}.
\end{problem}

\begin{problem}
已知椭圆 \(C: \frac{x^2}{a^2} + \frac{y^2}{b^2} = 1 (a > b > 0)\) 的焦距为 4，且过点 \(P(\sqrt{2}, \sqrt{3})\). (1) 求椭圆 C 的方程. (2) 设 \(Q(x_0, y_0) (x_0 \neq 0) \neq 0\) 为椭圆 \(C\) 上一点，过点 \(Q\) 作 \(x\) 轴的垂线，垂足为 \(E\). 取点 \(A(0, 2\sqrt{2})\)，连接 \(AE\). 过点 \(A\) 作 \(AE\) 的垂线交 \(x\) 轴于点 \(D\). 点 \(G\) 是点 \(D\) 关于 \(y\) 轴的对称点，作直线 \(QG\). 同样过作出的直线 \(QG\) 是否与椭圆 \(C\) 一定有唯一的公共点\(\perp\) 并说明理由 \fillinblank{}.
\end{problem}

